\documentclass[11pt]{article}
\usepackage[T1]{fontenc}
\usepackage{lmodern}
\usepackage{amsmath}
\usepackage{microtype}
\usepackage{graphicx}
\usepackage{booktabs}
\usepackage{xcolor}
\usepackage[numbers]{natbib}
\usepackage[colorlinks=true,linkcolor=blue,citecolor=blue!65!black,urlcolor=blue!65!black]{hyperref}
\usepackage[margin=1in]{geometry}

\title{Specialization in coupled learners: what worked and what failed}
\author{}
\date{}
\begin{document}
\maketitle

\section*{Question and model}

Starting from the diffusively coupled learners of
Arola-Fern\'andez and Lacasa \cite{arola2024}, we asked whether
\emph{dividing training effort} between learners could speed up the
learning of each individual network, and whether such a division
would emerge through adaptation.

We used four two-layer linear networks and two independent input
directions with a common identity teacher. All learners can solve
both tasks. A learner's allocation $a_i\in[0,1]$ sets its training
weight on the two task classes; its total local gradient budget
is fixed. Corresponding network parameters are diffusively
coupled with strength $\sigma$. The observable is mean
\emph{individual} test loss across both tasks, not ensemble accuracy.

\section*{What worked: specialization imposed in advance}

Fixed complementary specialists, with half the networks training
only on each task, can learn faster than uniformly trained
networks when parameter coupling transfers the missing knowledge.
At $\sigma=1$, initialization scale $0.03$, and a test-loss
threshold of $0.2$, median crossing times across 12 paired
seeds were $185$ SGD updates (uniform) and $150$ (specialists).
The specialists were faster in 11 of 12 pairs.

Without coupling, specialists did not reach this threshold
by 600 steps: learning only one task leaves the other unlearned.
At weaker coupling, $\sigma=0.2$, specialization was generally
slower. The advantage therefore requires an appropriate
combination of concentrated learning and knowledge transfer.

\section*{What failed: evolving the allocations}

We then allowed single learners to mutate their task allocations
at intervals of 20 or 60 updates. A proposed mutation was accepted
when a paired counterfactual rollout improved population-mean
validation loss. We compared it with uniform learning,
fixed specialists, frozen near-uniform allocations, and
neutral mutation drift, using 12 paired seeds per condition.

\begin{table}[htb!]
\centering
\begin{tabular}{lcc}
\toprule
Allocation & Steps to loss $0.2$ & Steps to loss $0.1$\\
\midrule
Uniform & 190 & 210\\
Fixed specialists & \textbf{163} & \textbf{195}\\
Frozen near-uniform & 188 & 210\\
Neutral drift & 190 & 218\\
Evolving & 193 & 228\\
\bottomrule
\end{tabular}
\caption{Mean first recorded crossing times for $\sigma=1$,
allocation updates every 20 steps, and 12 paired seeds;
all conditions reached both thresholds. Smaller is better.
Crossing times are resolved to 20-step intervals.}
\label{tab:outcomes}
\end{table}

The evolving allocations \emph{did not reproduce} the
learning-speed advantage of the fixed specialists.
They became heterogeneous, but at $\sigma=1$ their
final mean polarization was $0.406$, compared with $0.462$
under neutral drift. Differentiation alone therefore did
not demonstrate adaptive selection for specialization.

\begin{figure}[htb!]
\centering
\includegraphics[width=0.80\linewidth]{figures/arola_results.pdf}
\caption{Mean individual test loss over 12 seeds, with shaded
standard deviation; both panels use $\sigma=1$ and initialization
scale $0.03$. (a) Fixed allocations over 600 updates.
(b) Adaptive-allocation experiment over 400 updates, with
20-step selection windows. The dotted line marks the primary
test-loss threshold $0.2$.}
\label{fig:results}
\end{figure}

\section*{Why this is not yet an evolutionary theory}

The selection rule uses \emph{global}, held-out information and
extra counterfactual training. It does not describe individual
fitness, changing ecological competition, or decentralized
adaptation. Moreover, mutating one learner at a time changes
total effort on each task, even though complementary
specialization requires coordinated allocation.
The short validation horizon may also favor immediate
improvement over eventual learning speed. These are possible
explanations, not established causes of failure.

There is a separate analytical limitation. Allocation changes
both weighted input--output correlation and weighted input
covariance. Substituting an allocation-dependent coefficient
into the original scalar effective equation is therefore
\emph{not} a justified derivation.

\section*{Outcome}

The original effective-theory baseline was implemented;
the stored MNIST runs were short checks, not a full
reproduction of the published training protocol.
The microscopic experiment established that
\emph{imposed} specialization can accelerate learning,
but the attempted adaptive model did not explain how it
arises. We therefore retained these results and
shifted the main project to learned router--expert
dynamics in mixture-of-experts architectures, where
allocation and learning interact directly.

\bibliographystyle{unsrtnat}
\bibliography{refs}
\end{document}
